%% TeXora template: Accessible Media. Plum and cream 16:9 slides with a video and a GIF that play
%% inside the slide, each with what people who cannot see or hear it need: alt text, a caption, a
%% transcript (in the slides or in its own file) and no movement that starts and loops by itself.
%%
%% What it uses:
%%   - texoravideo (comes with TeXora; for another editor, keep texoravideo.sty and accessgate.sty
%%     beside this file). MP4 playback uses VLC or ffmpeg installed on your computer; GIFs need
%%     nothing.
%%   - media/sample-video.mp4, media/sample-loop.gif and media/sample-video-transcript.txt, which come
%%     with this template.
%%
%% The keys, in the second [...] of \includevideo:
%%   alt={...}            what the video is and what clicking does: "Play video: ..., 6 s"
%%   caption={...}        one short line under the video
%%   transcript={...}     the full text, printed where \AccTranscripts is, linked from the video
%%   transcriptfile=...   a transcript in its own file (.txt, .vtt or .pdf), linked from the video
%% Change the colours on the \definecolor lines below.
% !TEX program = pdflatex
\documentclass[aspectratio=169]{beamer}
\usetheme{default}
\usepackage[T1]{fontenc}
\usepackage{lmodern}
\usepackage{booktabs}
\usepackage{texoravideo}
\hypersetup{pdflang=en-US,pdftitle={Accessible Media},hidelinks}
\definecolor{plum}{RGB}{74,32,84}
\definecolor{cream}{RGB}{252,248,240}
\definecolor{rose}{RGB}{176,58,98}
\definecolor{blush}{RGB}{243,230,236}
\setbeamercolor{background canvas}{bg=cream}
\setbeamercolor{normal text}{fg=plum!85!black}
\setbeamercolor{structure}{fg=rose}
\setbeamercolor{frametitle}{bg=plum,fg=cream}
\setbeamercolor{title}{fg=plum}
\setbeamercolor{subtitle}{fg=rose}
\setbeamercolor{block title}{bg=rose,fg=white}
\setbeamercolor{block body}{bg=blush}
\setbeamertemplate{blocks}[rounded][shadow=false]
\setbeamertemplate{navigation symbols}{}
\setbeamertemplate{itemize items}[circle]
\setbeamertemplate{frametitle}{%
  \nointerlineskip\begin{beamercolorbox}[wd=\paperwidth,ht=3ex,dp=1.3ex,leftskip=1.2em]{frametitle}%
    \usebeamerfont{frametitle}\insertframetitle\end{beamercolorbox}}
\setbeamertemplate{footline}{%
  \hfill\usebeamercolor[fg]{structure}\scriptsize\insertframenumber\,/\,\inserttotalframenumber\hspace*{1.2em}\vskip0.8em}
\setbeameroption{show notes on second screen=right}

\newcommand{\code}[1]{{\small\ttfamily #1}}

\title{Accessible Media}
\subtitle{Videos and GIFs for everyone in the room}
\author{Your Name}
\institute{Your Institution}
\date{\today}

\begin{document}

\begin{frame}
  \vfill
  {\Huge\bfseries\usebeamercolor[fg]{title}\inserttitle\par}
  \vspace{0.5em}
  {\large\usebeamercolor[fg]{subtitle}\insertsubtitle\par}
  \vspace{1.6em}
  \insertauthor\quad\textperiodcentered\quad\insertinstitute\quad\textperiodcentered\quad\insertdate
  \vfill
\end{frame}
\note{Each clip in this deck has alt text and a text alternative. Check accessibility lists nothing to fix.}

\begin{frame}{A video with a caption}
  \begin{columns}[c]
    \begin{column}{0.52\textwidth}
      \includevideo[width=\linewidth][posterat=30,alt={Play video: a light beam turning, 6 seconds, no sound},caption={A light beam turns once around the centre.}]{media/sample-video.mp4}
    \end{column}
    \begin{column}{0.44\textwidth}
      \code{[alt=\{Play video: a light beam}\\
      \code{\ \ turning, 6 seconds, no sound\},}\\
      \code{\ caption=\{A light beam turns ...\}]}
      \par\medskip
      \begin{itemize}
        \item The poster is a link, so \textbf{alt} names what clicking does.
        \item The \textbf{caption} shows under the video, for everyone.
      \end{itemize}
    \end{column}
  \end{columns}
\end{frame}
\note{Click the clip to play it. Say the caption out loud too: it helps people at the back.}

\begin{frame}{A transcript in the slides}
  \begin{columns}[c]
    \begin{column}{0.52\textwidth}
      \includevideo[width=\linewidth][alt={Play video: a light beam turning, 6 seconds},transcript={The screen is dark. A single beam of light appears at the left and turns clockwise around the centre, once, in six seconds. There is no speech and no music.}]{media/sample-video.mp4}
    \end{column}
    \begin{column}{0.44\textwidth}
      \begin{itemize}
        \item \code{transcript=\{...\}}: any length.
        \item A \emph{Transcript} link under the video leads to the text, printed on the last slide by \code{\textbackslash AccTranscripts}.
        \item Describe what is seen as well as what is said.
      \end{itemize}
    \end{column}
  \end{columns}
\end{frame}

\begin{frame}{A transcript in its own file}
  \begin{columns}[c]
    \begin{column}{0.52\textwidth}
      \includevideo[width=\linewidth][alt={Play video: a light beam turning, 6 seconds},transcriptfile=media/sample-video-transcript.txt]{media/sample-video.mp4}
    \end{column}
    \begin{column}{0.44\textwidth}
      \begin{itemize}
        \item \code{transcriptfile=media/...txt}: a \code{.txt}, \code{.vtt} (subtitles) or \code{.pdf} kept with the slides.
        \item The link under the video opens it.
        \item Good for long talks and interviews.
      \end{itemize}
    \end{column}
  \end{columns}
\end{frame}

\begin{frame}{Movement you can stop}
  \begin{columns}[c]
    \begin{column}{0.48\textwidth}
      \centering
      \includevideo[height=4cm][loop,alt={Play animation: dots drifting from left to right, it loops},caption={Dots drift across the picture. Click to start or stop.}]{media/sample-loop.gif}
    \end{column}
    \begin{column}{0.48\textwidth}
      \begin{itemize}
        \item Movement that starts by itself, loops and lasts over 5 seconds distracts many people (WCAG 2.2.2).
        \item Here the GIF \textbf{loops} but starts on a click, and a click stops it.
        \item The check flags \code{autoplay} together with \code{loop}.
      \end{itemize}
    \end{column}
  \end{columns}
\end{frame}
\note{If a looping clip must start by itself, keep it short, or tell the audience that a click pauses it.}

\begin{frame}{Which key for which clip}
  \centering
  \AccHeaderRows{1}
  \begin{tabular}{@{}lll@{}}
    \toprule
    \textbf{Clip} & \textbf{Needs} & \textbf{Key} \\
    \midrule
    Silent, short & what it shows & \code{caption} \\
    Speech or sound & every word, and the sounds & \code{transcript} or \code{transcriptfile} \\
    Long, with subtitles & the subtitle file & \code{transcriptfile=talk.vtt} \\
    Decorative loop & nothing to read & \code{alt}, and no \code{autoplay} \\
    \bottomrule
  \end{tabular}
\end{frame}

\begin{frame}[allowframebreaks]{Transcripts}
  \AccTranscripts
\end{frame}
\note{The transcripts frame breaks over as many slides as it needs.}

\end{document}
